% Part: first-order-logic
% Chapter: sequent-calculus
% Section: rules-and-proofs

\documentclass[../../../include/open-logic-section]{subfiles}

\begin{document}

\iftag{FOL}
      {\olfileid{fol}{seq}{rul}}
      {\olfileid{pl}{seq}{rul}}

\olsection{नियम आणि \usetoken{P}{derivation}}

पुढील मजकुरात, $\Gamma, \Delta, \Pi, \Lambda$ या !!{sentence}s च्या सांत
क्रमिका दर्शवतील, असे मानू.

\begin{defn}[क्रमवर्ती]
\emph{क्रमवर्ती} म्हणजे पुढील रूपाची अभिव्यक्ती होय:
\[
\Gamma \Sequent \Delta
\]
येथे $\Gamma$ आणि $\Delta$ या $\Lang L$ भाषेतील !!{sentence}s च्या
सांत (संभवतः रिक्त) क्रमिका आहेत. $\Gamma$ ला \emph{पूर्वांग}, तर
$\Delta$ ला \emph{उत्तरांग} म्हणतात.
\end{defn}

\begin{explain}
क्रमवर्तीमागील अंतःप्रेरित कल्पना अशी आहे: पूर्वांगातील सर्व !!{sentence}s
सत्य असतील, तर उत्तरांगातील !!{sentence}s पैकी किमान एक सत्य असते. म्हणजे,
$\Gamma = \tuple{!A_1, \dots, !A_m}$ आणि
$\Delta = \tuple{!B_1, \dots, !B_n}$ असतील, तर $\Gamma \Sequent \Delta$
तेव्हा आणि केवळ तेव्हाच सत्य असते, जेव्हा
\[
(!A_1 \land \cdots \land !A_m) \lif (!B_1 \lor \cdots \lor
!B_n)
\]
सत्य असते. याची दोन विशेष प्रकरणे आहेत: $\Gamma$ रिक्त असण्याचे आणि
$\Delta$~रिक्त असण्याचे. $\Gamma$ रिक्त असेल, म्हणजेच $m = 0$ असेल, तर
$\quad \Sequent \Delta$ तेव्हा आणि केवळ तेव्हाच सत्य असते, जेव्हा
$!B_1 \lor \dots \lor !B_n$ सत्य असते. $\Delta$ रिक्त असेल, म्हणजेच
$n = 0$ असेल, तर $\Gamma \Sequent \quad$ तेव्हा आणि केवळ तेव्हाच सत्य
असते, जेव्हा $\lnot(!A_1 \land \dots \land !A_m)$ सत्य असते. एखाद्या
क्रमवर्तीला तेव्हा आणि केवळ तेव्हाच वैध म्हणतो, जेव्हा तिच्याशी संबंधित
!!{sentence} वैध असते.
\end{explain}

$\Gamma$ ही !!{sentence}s ची क्रमिका असेल, तर $\Gamma$ च्या उजव्या टोकाला
$!A$ जोडून मिळणाऱ्या क्रमिकेसाठी आपण $\Gamma, !A$ लिहितो (आणि $\Gamma$
च्या डाव्या टोकाला $!A$ जोडून मिळणाऱ्या क्रमिकेसाठी $!A, \Gamma$
लिहितो). $\Delta$ हीदेखील !!{sentence}s ची क्रमिका असेल, तर $\Gamma,
\Delta$ ही त्या दोन क्रमिकांची जोडणी असते.

\begin{defn}[आरंभीची क्रमवर्ती]
\emph{आरंभीची क्रमवर्ती} म्हणजे, भाषेतील कोणत्याही !!{sentence}~$!A$ साठी,
\iftag{prvFalse,prvTrue}{पुढीलपैकी एका रूपाची क्रमवर्ती होय:
  \begin{enumerate}
    \item $!A \Sequent !A$
    \tagitem{prvTrue}{$\quad \Sequent \ltrue$}{}
    \tagitem{prvFalse}{$\lfalse \Sequent \quad$}{}
  \end{enumerate}}
{$!A \Sequent !A$ या रूपाची क्रमवर्ती होय}.
\end{defn}

क्रमवर्ती कलनातील !!^{derivation}s हे क्रमवर्तींचे विशिष्ट वृक्ष असतात.
त्यांतील सर्वांत वरच्या क्रमवर्ती आरंभीच्या क्रमवर्ती असतात; आणि एखादी
क्रमवर्ती इतर एक किंवा दोन क्रमवर्तींच्या खाली असेल, तर ती निगमनाच्या
एखाद्या नियमानुसार त्यांच्यापासून योग्य रीतीने निष्पन्न झाली पाहिजे.
$\Log{LK}$ चे नियम दोन मुख्य प्रकारांत विभागले आहेत: \emph{तार्किक} नियम
आणि \emph{संरचनात्मक} नियम. खालच्या क्रमवर्तीतील $!A$ आणि/किंवा $!B$
असलेल्या !!{sentence} च्या !!{main
  operator} वरून तार्किक नियमांना नावे दिली
जातात. प्रत्येक नियमाच्या दोन आवृत्त्या असतात: !!{operator} असलेले
!!{sentence} डावीकडे असणारी क्रमवर्ती निष्पन्न करण्यासाठी एक, आणि ते
!!{sentence} उजवीकडे असणारी क्रमवर्ती निष्पन्न करण्यासाठी दुसरी.

\end{document}
